\documentclass{article}
\usepackage[utf8]{inputenc}
\usepackage[T1]{fontenc}
\usepackage{geometry}
\geometry{
    a4paper,
    left=2.5cm,
    right=2.5cm,
    top=2.5cm,
    bottom=2.5cm
}
\usepackage{graphicx}
\usepackage{float}
\usepackage{hyperref}
\usepackage[english]{babel}
\usepackage{amsmath}
\usepackage{amssymb}
\usepackage{array}
\usepackage{booktabs}
\usepackage{caption}
\usepackage{subcaption}
\usepackage{enumitem}
\usepackage{fancyhdr}
\usepackage{titlesec}
\usepackage{url}
\usepackage{textcomp}
\usepackage{xcolor}
\usepackage{listings}
\usepackage{multirow}
\usepackage{longtable}
\usepackage{tabularx}
\usepackage{times}

% Set up page style
\pagestyle{fancy}
\fancyhf{}
\rfoot{\thepage}
\renewcommand{\headrulewidth}{0pt}

% Title formatting
\titleformat{\chapter}[display]{\normalfont\bfseries}{\chaptertitle\thechapter}{0pt}{\Huge}
\titlespacing*{\chapter}{0pt}{50pt}{40pt}

% Document info
\title{Template Document}
\author{Academic Support Agent}
\date{\today}

\begin{document}
\maketitle

\pagenumbering{roman} %Saavuttavuuden takija alussa sivunurot roomalaisittain, tutkielman alusta arabialaisittain.
\pagestyle{empty}  %ei sivunumeroa sivun alareunaan

\begin{center}
\includegraphics[width=10cm]{UTU_logo_FI} %talleta kuva linkist  omaan hakemistoosi
\end{center}

\vspace{3.0cm}
\begin{center}\large
{\sc \otsikko} 
\end{center}

\vspace{0.5cm}
\begin{center}
\titteli \tekija
\end{center}

\vspace{0.5cm}
\begin{center}
\tutkielma -tutkielma\\
\aika
\end{center}

\vspace{2.5cm}
\begin{center}
\begin{tabular}{l}
Tarkastajat:\\
\ohjaaja \\
\tarkastaja
\end{tabular}
\end{center}


\vspace{2.5cm}
\begin{center}
MATEMATIIKAN JA TILASTOTIETEEN LAITOS
\end{center}

\newpage\null

\vspace{22cm}

%Pakollinen ilmoitus Turnitin-j rjestelm n k yt st :
\noindent Turun yliopiston laatuj rjestelm n mukaisesti t m n julkaisun alkuper isyys on tarkastettu
Turnitin OriginalityCheck-j rjestelm ll 






\noindent
TURUN YLIOPISTO\\
Matematiikan ja tilastotieteen laitos\\
\\
{\sc \tekija}: \otsikko \\
\tutkielma -tutkielma, 39 s., 12 liites.\\ %t yt  sivum r t!
\paaaine \\ 
\aika \\
\rule{\textwidth}{.2mm}\\
\\

\vspace{4mm}\noindent Kirjoita t h n tiivistelm . Laita $\backslash${noindent}-komento kappaleen
alkuun, niin \LaTeX\, ei sisenn  ensimm ist  rivi .

\vspace{4mm}\noindent Komento $\backslash$vspace taasen j tt  sopivan v lin kappaleiden v liin - 4mm n ytt  aika hyv lt .

\vspace{4mm}\noindent Kirjoita tiivistelm  napakasti ja kaikenlaista toistoa v ltt en.

\vspace{4mm}\noindent Asiasanat: tiivistelm sivu, Pro gradu -tutkielma, \LaTeX-ladontaj rjestelm .







%Aja laTeX-k nn s uudelleen saadaksesi ajantasaisen sis llysluettelon



\pagestyle{plain} 
\pagenumbering{arabic} % Saavutettavuus: Aloitetaan arabialaisilla numeroilla t lt  sivulta

\section{Integraalilaskentaa}

T ss  luvussa esitet n niin kutsuttu Newtonin-Leibnizin kaava \cite{NewtLeib}.

\begin{theorem} Oletetaan ett  v lill  $[a,b]$ on $F'(x)=f(x)$ ja ett  $f$ on jatkuva.\footnote{Jatkuvuus on oleellinen seikka t ss  yhteydess .} T ll in
\[
\int_a^bf(x)\,dx=F(b)-F(a).
\]
\end{theorem}

\subsection{Analyyttist  lukuteoriaa}

Seuraavaksi esitell n niin sanottu {\em Riemannin hypoteesi}, jonka mukaan ns. $\zeta$-funktion ep triviaalit nollakohdat sijaitsevat suoralla $\operatorname{Re}z=\frac12$.
\begin{definition} Riemannin $\zeta$-funktio m ritell n sarjaesityksen \cite{Riemann}
\[
\zeta(z)=\sum_{n=1}^{\infty}\frac{1}{n^z}
\]
avulla, kun $\operatorname{Re}z>1$.
\end{definition}

\subsubsection{Matriisilaskentaa}

Nyt tarkastellaan matriisin
\[
A=\left(\begin{array}{rrr} -1 & -2 & -5\\ 3& 4& 5\\ -3 & 2 & 1 \end{array}\right)
\]
ominaisarvoja.

\section{Saavutettavuus ja kuvat}

UTUGradu-j rjestelm n ja PDF\LaTeX :n pdf-formaattivaatimusten takia kannattaa k ytt  muita kuin eps-muotoisia kuvia. Esimerkiksi pdf-muoto k y mainiosti vektorikuville, kunhan on samaa pdf/a-3b muotoa kuin t m n tiedoston tuottama pdf-tiedosto. Parhaiten toimivat jpg-muotoiset kuvat.



\subsection{Saavutettavuus kuvateksteiss }

Kaikkiin tutkielmassa esiintyviin kuviin tulee viitata tekstiss  (ks.~\ref{kuvatus1}) Lis ksi kuvatekstin tulee olla kuvaileva, koska saavutettavuuteen liittyv t alt-tekstit eiv t oikein toimi k ytetyn \LaTeX :n axessibility-paketin kanssa. Jos kuvan perusteella tehd n p telm , tulee p telm  kuvata tekstiss . 

\begin{figure}[ht]
\begin{center}
\includegraphics[width=.45\textwidth]{siivet.jpg}
\end{center}
\label{kuvatus1}
\caption{Geronon lemniskaatta v lill  $[-1,1]$.}
\end{figure}

\section{Teksti }

\lipsum[1-10]


%Kirjallisuusluettelon m rittelyss  {99} on levein kirjallisuusviitteen numero. Oikaise tarvittaessa.
 
\begin{thebibliography}{99}

\bibitem{NewtLeib} Newton ja Leibniz

\bibitem{Riemann} Riemann

\end{thebibliography}
\end{document}